\documentclass[11pt, a4paper]{article}

\usepackage[utf8]{inputenc}
\usepackage{geometry}
\geometry{a4paper, margin=0.5in}
\usepackage{xcolor}
\usepackage{titlesec}
\usepackage[listings,breakable,skins]{tcolorbox}
\usepackage{listings}
\pagenumbering{gobble}

\definecolor{darkgray}{RGB}{40, 40, 40}
\definecolor{promptbg}{RGB}{245, 245, 250}
\definecolor{promptborder}{RGB}{100, 100, 200}

\titleformat{\section}{\large\bfseries\color{darkgray}}{\thesection}{1em}{}[\titlerule]
\titleformat{\subsection}{\normalsize\bfseries\color{darkgray}}{\thesubsection}{1em}{}

\lstdefinestyle{promptstyle}{
    basicstyle=\small\ttfamily,
    breaklines=true,
    breakatwhitespace=false,
    columns=fullflexible,
    keepspaces=true,
    showstringspaces=false,
    upquote=true,
    literate={`}{{\textasciigrave}}1
}

\newtcblisting{promptbox}{
    listing only,
    breakable,
    enhanced,
    colback=promptbg,
    colframe=promptborder,
    boxrule=0.5pt,
    arc=4pt,
    left=6pt, right=6pt, top=6pt, bottom=6pt,
    listing options={style=promptstyle}
}


\title{\textbf{Prompt Engineering Playbook v2} \\ \large Auxin Automata --- LLM Architecture \& Code Generation}
\author{Edwin Redhead \& Tara Kasayapanand}
\date{Colosseum Frontier Hackathon Sprint}

\begin{document}
\maketitle

\begin{abstract}
Master prompts for generating the Auxin Automata MVP infrastructure. Each prompt is paired with the highest-performing model and IDE for the task. Prompts are written to enforce the SDK-first, agnosticism-by-code-shape architecture defined in the Technical Roadmap v2.
\end{abstract}

\vspace{1em}

\section{Phase 0: Repository Bootstrap}

\subsection*{0.1--0.6: Monorepo, CI, Devcontainer}
\textbf{Tool:} Claude Code (terminal, repo root) \quad \textbf{Model:} Claude Sonnet 4.6

\begin{promptbox}
Context: Bootstrapping the auxin-automata monorepo for a 5-week Solana hackathon (Colosseum Frontier). Two devs, parallel tracks. Python SDK + Anchor programs + Next.js dashboard + ROS2 edge.

Task: Create the full repo skeleton.

Requirements:
1. Workspaces: /sdk (Python, uv-managed), /programs (Anchor), /edge (ROS2 Python nodes), /dashboard (Next.js 14 app router + TS strict + Tailwind), /twin (PyBullet), /docs, /scripts.
2. Root files: README.md (architecture placeholder), LICENSE (Apache 2.0), .gitignore (Python + Node + Rust + Anchor + .env), CODEOWNERS, .editorconfig, Makefile with targets: bootstrap, lint, test, demo.
3. .github/workflows/: ci.yml with parallel jobs for python-lint-test (ruff + pytest), anchor-build-test, dashboard-build-lint. Fail-fast matrix.
4. .devcontainer/devcontainer.json pinning Python 3.11, Node 20, Rust stable, Solana CLI 1.18+, Anchor 0.30+.
5. Pre-commit config with gitleaks, ruff, prettier.
6. .env.example in each workspace with documented variables (no real secrets).
7. PR template and issue templates (bug, feature) under .github/.

Output as a tree first, then generate every file. Do not skip files.
\end{promptbox}

\subsection*{0.7: Competitive Scan (Colosseum Copilot)}
\textbf{Tool:} Claude Code with Colosseum Copilot skill installed \quad \textbf{Model:} Claude Sonnet 4.6

\begin{promptbox}
Vet this idea: Auxin Automata is an SDK + Solana program suite that gives physical hardware (robotic arms, drones, industrial machines) its own on-chain wallet, lets it autonomously stream M2M micropayments for AI inference and external compute, and hashes kinematic safety telemetry to Solana as an immutable compliance log. MVP target: a desktop robotic arm running ROS2 on a Jetson, paying Gemini for vision inference, and logging torque-anomaly events on Devnet.

Run a deep dive. Specifically:
1. Have any prior Colosseum hackathon teams (Renaissance, Radar, Breakout, Cypherpunk) shipped anything in the autonomous-hardware-wallet, M2M payments for compute, or on-chain robotic compliance space? Name them.
2. Classify the gap (full / partial / false). If partial, on which axis (segment, UX, geographic, integration)?
3. Funded competitors in DePIN + agentic payments + machine identity. Names, raise amounts, latest shipped product.
4. Recommend a wedge: what's the most defensible narrow market we should anchor the pitch around (surgical robotics audit trails? warehouse automation insurance? drone airspace logging?).
5. Cite every claim. If evidence floor isn't met, say so explicitly.

Output as a structured report we can paste into the pitch deck appendix.
\end{promptbox}

\section{Phase 1: Foundation \& Mock Infrastructure}

\subsection*{1A: SDK Core (\texttt{auxin-sdk})}
\textbf{Tool:} Cursor IDE \quad \textbf{Model:} Claude Sonnet 4.6

\begin{promptbox}
Context: Building auxin-sdk, the Python package that proves our hardware-agnostic architecture to Colosseum judges. The same SDK must drive a mock generator, a PyBullet digital twin, and a real ROS2 robot --- selected by an env var.

Task: Scaffold the auxin-sdk package under /sdk.

Requirements:
1. pyproject.toml with uv, src/auxin_sdk/ layout, Python 3.11+. Deps: solders, solana, pydantic v2, structlog, anyio, httpx, google-generativeai, tenacity.
2. auxin_sdk/wallet.py: HardwareWallet class wrapping solders Keypair. Methods: load_or_create(path), pubkey, sign_transaction(tx), get_balance(rpc_url), request_airdrop(rpc_url, sol). Async where it touches the network.
3. auxin_sdk/schema.py: Pydantic v2 TelemetryFrame model with fields timestamp (datetime), joint_positions (list[float]), joint_velocities (list[float]), joint_torques (list[float]), end_effector_pose (dict), anomaly_flags (list[str]). Include model_config for JSON serialisation.
4. auxin_sdk/hashing.py: canonical_json(frame) using sorted keys and fixed float repr; sha256_hex(frame). Include a determinism test fixture.
5. auxin_sdk/sources/base.py: TelemetrySource ABC with `async def stream() -> AsyncIterator[TelemetryFrame]` and `async def close()`.
6. auxin_sdk/logging.py: configure_structlog() emitting JSON with request_id contextvar.
7. tests/ with pytest covering wallet round-trip, schema validation, hash determinism. Target 80% coverage.
8. Add SDK to root Makefile lint/test targets.

Do NOT implement concrete sources here --- they go in 1B/1C/2B. Just the contract.
\end{promptbox}

\subsection*{1B: Mock Telemetry Source}
\textbf{Tool:} Cursor IDE \quad \textbf{Model:} Claude Sonnet 4.6

\begin{promptbox}
Context: Extending auxin-sdk with a mock telemetry source that conforms to the TelemetrySource ABC defined in /sdk/src/auxin_sdk/sources/base.py.

Task: Implement auxin_sdk/sources/mock.py.

Requirements:
1. MockSource(TelemetrySource) class. Constructor accepts: rate_hz (default 10), num_joints (default 6), anomaly_every (default 12, randomised +/-3), seed (optional int for reproducibility).
2. stream() yields TelemetryFrame objects forever. Joint positions drift via per-joint sin/cos with phase offset + small Gaussian noise. Velocities are numerical derivative. Torques baseline + noise.
3. Every anomaly_every frames, spike one random joint's torque to 95.0 and append "torque_spike" to anomaly_flags.
4. Add a JSONL recorder (record_to(path)) and a ReplaySource(path) class for deterministic demo replays.
5. Curate /sdk/fixtures/images/ with 20 placeholder image filenames (clear_01.jpg .. obstacle_10.jpg) and a labels.json mapping filename -> "clear"|"obstacle". Add a sampler function sample_workspace_image() returning (path, label).
6. Tests: assert frames validate against schema, assert anomalies appear within expected window, assert ReplaySource reproduces a recorded session bit-identical.
\end{promptbox}

\subsection*{1C: Digital Twin Fallback}
\textbf{Tool:} Cursor IDE \quad \textbf{Model:} Claude Opus 4.6 (PyBullet specifics need deeper reasoning)

\begin{promptbox}
Context: Building the digital twin fallback for Auxin Automata. This is the insurance policy: if the physical robotic arm slips, the twin must carry the entire demo video. It must conform to the same TelemetrySource ABC as MockSource so the bridge code is identical.

Task: Build /twin as a standalone module that exposes a TwinSource.

Requirements:
1. /twin/scene.py: PyBullet scene loader. Load Franka Panda URDF from pybullet_data, a table, and one obstacle box. Provide reset() and step() at 240 Hz internal sim rate.
2. /twin/trajectory.py: Pre-scripted pick-and-place loop using PyBullet's calculateInverseKinematics. Loops indefinitely.
3. /twin/source.py: TwinSource(TelemetrySource) class subclassing the SDK ABC. Internally runs the sim in a background task; stream() yields TelemetryFrames at 10 Hz built from getJointStates. Map PyBullet joint state to the schema fields exactly.
4. /twin/render.py: Off-screen renderer producing (a) MP4 file via imageio + ffmpeg AND (b) an asyncio websocket server on ws://localhost:8765 streaming JPEG frames base64-encoded for the dashboard.
5. /twin/__main__.py: CLI entrypoint `python -m twin --mode [video|ws|both]`.
6. The bridge selects this source via AUXIN_SOURCE=twin --- do NOT special-case it in any downstream code.
7. Tests: smoke test that 100 frames stream cleanly and validate against schema.

Critical: keep the TwinSource interface byte-identical to MockSource and (eventually) ROS2Source. Agnosticism must show up in `git diff` as a one-line config change.
\end{promptbox}

\subsection*{1D: Gemini Safety Oracle}
\textbf{Tool:} Cursor IDE \quad \textbf{Model:} Claude Sonnet 4.6 (test the system prompt itself in Gemini's web UI first)

\begin{promptbox}
Context: Adding the Gemini-backed safety oracle to auxin-sdk. This is the AI decision layer that approves or denies robotic actions based on telemetry + workspace image.

Task: Implement auxin_sdk/oracle.py.

Requirements:
1. SafetyOracle class. Constructor: api_key (from env GEMINI_API_KEY), model (default "gemini-2.0-flash"), timeout_s (default 2.0), torque_threshold (default 80.0).
2. async def check(frame: TelemetryFrame, image_path: Path) -> OracleDecision. OracleDecision is a Pydantic model: action_approved (bool), reason (str), confidence (float), latency_ms (float), prompt_version (str), used_fallback (bool).
3. Use Gemini's response_schema feature to enforce structured JSON output. Schema must match OracleDecision exactly minus latency_ms/used_fallback.
4. System prompt loaded from /sdk/src/auxin_sdk/prompts/safety_oracle_v1.txt. Prompt instructs Gemini to act as an industrial safety oracle: deny if image shows an obstacle, deny if any torque > threshold, otherwise approve. Log the prompt version used.
5. Retry with tenacity: 3 attempts, exponential backoff, only on transient errors.
6. On total failure or timeout, fall back to local heuristic: deny if max(torques) > threshold OR image label (from fixtures/labels.json if available) == "obstacle". Set used_fallback=True. The demo must NEVER stall on a network blip.
7. Cost/latency metrics: record token counts and latency to a structlog event "oracle.decision".
8. Eval harness in tests/eval_oracle.py: run the oracle against all 20 fixture images with synthetic safe + unsafe telemetry, assert >=90% accuracy. CI runs this nightly only (it costs API credits).
\end{promptbox}

\section{Phase 2A: Blockchain (Solana / Anchor)}

\subsection*{2A.1--2A.2: Anchor Workspace + PDA Design}
\textbf{Tool:} Claude Code (terminal) \quad \textbf{Model:} Claude Sonnet 4.6

\begin{promptbox}
Context: Scaffolding the Anchor workspace for Auxin Automata's M2M payment + compliance bridge on Solana Devnet. This program is consumed by a Python SDK (auxin-sdk) and a Next.js dashboard.

Task: Create /programs/agentic_hardware_bridge as an Anchor 0.30+ workspace.

Step 1 --- Design doc: Before writing Rust, output /programs/DESIGN.md describing the three PDAs:
- HardwareAgent: seeds=[b"agent", owner_pubkey], fields: owner, hardware_pubkey, compute_budget_lamports, lamports_spent, providers (Vec<Pubkey>, max 8), created_at, bump.
- ComputeProvider: seeds=[b"provider", provider_pubkey], fields: provider_pubkey, total_received, bump.
- ComplianceLog: seeds=[b"log", agent_pubkey, slot.to_le_bytes()], fields: agent, hash (String, max 64), severity (u8), reason_code (u16), timestamp, bump.

Step 2 --- Scaffold: anchor init agentic_hardware_bridge. Configure Anchor.toml for Devnet via env HELIUS_RPC_URL. lib.rs imports + declare_id placeholder. Modules: state.rs (PDA structs), errors.rs, instructions/ (one file per ix), events.rs.

Step 3 --- Test harness: Set up TS tests with @coral-xyz/anchor, helper to create funded keypairs on localnet, helper to derive each PDA.

Do NOT implement instruction logic yet --- just scaffolding, types, and the design doc. We'll fill instructions in 2A.3--2A.6.
\end{promptbox}

\subsection*{2A.3--2A.7: Instructions, Events, Errors}
\textbf{Tool:} Claude Code \quad \textbf{Model:} Claude Opus 4.6

\begin{promptbox}
Context: Implementing the instruction set for agentic_hardware_bridge per /programs/DESIGN.md.

Task: Write all four instructions, the events, and the error enum.

Requirements:
1. errors.rs: #[error_code] enum AuxinError with variants UnauthorizedSigner, BudgetExceeded, RateLimitExceeded, InvalidProvider, ProviderNotWhitelisted, HashTooLong, MaxProvidersReached.

2. events.rs: ComputePaymentEvent { agent, provider, lamports, timestamp }, ComplianceEvent { agent, hash, severity, reason_code, timestamp }, AgentInitializedEvent, ProviderWhitelistUpdatedEvent.

3. instructions/initialize_agent.rs: Creates HardwareAgent PDA. Args: hardware_pubkey, compute_budget_lamports. Constraints: owner is signer, init account.

4. instructions/stream_compute_payment.rs: Args: lamports. Constraints: hardware_pubkey is signer (autonomous --- the hardware itself signs, not the owner), provider PDA in agent.providers. Logic: enforce per-tx cap (max 0.001 SOL), enforce rolling-window rate limit (max 100 tx per 60 slots --- store last_window_start_slot and window_count on the agent), check budget remaining, transfer lamports via system_program::transfer with PDA signer seeds, update lamports_spent, emit ComputePaymentEvent.

5. instructions/log_compliance_event.rs: Args: hash (String, max 64 chars --- enforce), severity (u8, 0--3), reason_code (u16). Constraints: hardware_pubkey is signer. Logic: init ComplianceLog PDA, store fields, emit ComplianceEvent. Compliance logs can never be rate-limited or budget-blocked.

6. instructions/update_provider_whitelist.rs: Args: provider, action (Add|Remove enum). Constraints: owner is signer. Enforce max 8 providers.

7. TypeScript test suite covering: happy path for each ix, every error variant, multi-agent isolation (agent A cannot pay from agent B's budget), event listener verifying every emitted event field, rate limit boundary (101st tx in window fails, then passes after window rolls).

Use the latest Anchor macro syntax. No deprecated patterns.
\end{promptbox}

\subsection*{2A.10--2A.12: Devnet Deploy, Client Gen, Security}
\textbf{Tool:} Claude Code \quad \textbf{Model:} Claude Sonnet 4.6

\begin{promptbox}
Context: agentic_hardware_bridge passes all localnet tests. Now we deploy to Devnet, generate clients, and run a security pass.

Task:
1. /scripts/deploy_devnet.sh: idempotent deploy. Reads HELIUS_RPC_URL and DEPLOYER_KEYPAIR_PATH from env. Builds with `anchor build`, deploys with `anchor deploy --provider.cluster devnet`, uploads IDL, writes program ID to /programs/deployed.json with cluster + timestamp. Re-running must not redeploy unless --force is passed.

2. Post-deploy smoke test script (TypeScript): initialize a test agent, stream one payment, log one compliance event, assert all events fire. Outputs Explorer links.

3. Python client generation: use anchorpy to generate /sdk/src/auxin_sdk/program/ from the IDL. Wrap it in a high-level AuxinProgramClient class with methods initialize_agent, stream_payment, log_compliance --- each accepting a HardwareWallet and returning the tx signature.

4. Security pass checklist (output as /programs/SECURITY.md):
   - Every instruction validates the correct signer (hardware vs owner).
   - No CPI to untrusted programs.
   - Integer arithmetic uses checked_add / checked_sub.
   - All account constraints use has_one or seeds where appropriate.
   - String length bounds enforced before storage.
   - Run cargo-audit and paste results.
   - Run cargo clippy --all-targets -- -D warnings.
\end{promptbox}

\section{Phase 2B: Edge Compute (Jetson / ROS2)}

\subsection*{2B.5--2B.6: Telemetry Bridge \& Watchdog}
\textbf{Tool:} Cursor IDE over SSH to Jetson \quad \textbf{Model:} Claude Opus 4.6

\begin{promptbox}
Context: ROS2 Humble on NVIDIA Jetson Orin Nano. A myCobot 280 publishes /joint_states. We need a bridge node that conforms to auxin-sdk's TelemetrySource ABC, plus an INDEPENDENT watchdog node that can halt the arm even if the network or Solana bridge dies.

Task: Build two ROS2 Python nodes under /edge/auxin_edge/.

Node 1 --- telemetry_bridge_node.py:
1. Subscribe to /joint_states (sensor_msgs/JointState) with QoS BEST_EFFORT.
2. Throttle: keep only the latest message, emit at 2 Hz via a timer.
3. Convert to auxin_sdk.schema.TelemetryFrame (import from the installed SDK package).
4. Push to an asyncio queue exposed by a ROS2Source(TelemetrySource) class also defined in this file. The bridge process imports ROS2Source and uses it identically to MockSource/TwinSource.
5. Handle dropped messages: if no joint_states for >1s, log warning, push a frame with anomaly_flag "stale_telemetry".
6. Graceful shutdown on SIGINT.

Node 2 --- safety_watchdog_node.py:
1. Subscribe to /joint_states at full rate.
2. If max(abs(effort)) > 80.0 for 3 consecutive messages, call the /emergency_stop service (std_srvs/Trigger).
3. Publishes its own /auxin/watchdog_status topic (heartbeat).
4. Runs as a SEPARATE process under systemd. Must not import auxin-sdk, must not touch the network. Independence is the safety guarantee.

Add /edge/launch/auxin_edge.launch.py launching both nodes, and /edge/systemd/ unit files for boot autostart.
\end{promptbox}

\section{Phase 2C: Dashboard (Next.js)}

\subsection*{2C.1--2C.5: Mechafloral UI Scaffold}
\textbf{Tool:} Cursor IDE (Composer) \quad \textbf{Model:} Claude Sonnet 4.6

\begin{promptbox}
Context: Building the Auxin Automata dashboard --- the visual proxy judges use to understand the backend. Aesthetic is "Mechafloral": dark slate gray base, deep teal + Solana green accents, geometric clean layout, highly technical and venture-funded looking.

Task: Scaffold the dashboard at /dashboard with mock data first. Live RPC hookup comes in a separate prompt (2C.6).

Stack: Next.js 14 app router, TypeScript strict, Tailwind, shadcn/ui, lucide-react, Recharts, Zustand.

Requirements:
1. tailwind.config.ts: define design tokens. Background #0a0e1a, surface #131826, border #1f2937, accent-teal #14b8a6, accent-solana #14F195, text-primary #f1f5f9, text-muted #94a3b8. Inter font.
2. app/layout.tsx: dark mode locked, Inter loaded, root container.
3. app/page.tsx: header + three-column responsive grid (telemetry left, payments + compliance right column stacked, twin viewport spanning).
4. components/Header.tsx: "Auxin Automata --- Agentic Infrastructure Node" title, truncated agent pubkey with copy button, pulsing green "Status: Live" indicator (lucide Activity icon).
5. components/TelemetryCard.tsx: live joint angles + torques. Recharts sparklines per joint. Card border shifts to red on anomaly.
6. components/PaymentTicker.tsx: terminal-style scrolling list. Each row: timestamp, lamports, provider pubkey (truncated), tx signature linking to explorer.solana.com/?cluster=devnet.
7. components/ComplianceTable.tsx: shadcn Table. Columns: timestamp, severity badge (color-coded 0--3), reason code, hash (truncated, copy on click), tx link.
8. components/TwinViewport.tsx: react-three-fiber canvas, OrbitControls, placeholder cube for now (real URDF in a follow-up).
9. lib/store.ts: Zustand store with telemetry, payments, complianceLogs arrays + setters. Mock data in lib/mockData.ts feeds all four components on a setInterval.
10. Every component handles loading, empty, and error states. No white screens.
11. Vercel-ready: no env vars required for the mock build.
\end{promptbox}

\subsection*{2C.6: Live Solana Hookup (\texttt{useProgramEvents})}
\textbf{Tool:} Cursor IDE \quad \textbf{Model:} Claude Sonnet 4.6

\begin{promptbox}
Context: The dashboard mock UI is built. Now wire it to the deployed agentic_hardware_bridge program on Devnet, replacing mock data with real on-chain events.

Task: Build a useProgramEvents hook and swap mock data for live data.

Requirements:
1. lib/solana.ts: Connection factory reading NEXT_PUBLIC_HELIUS_RPC_URL and NEXT_PUBLIC_PROGRAM_ID. Export PROGRAM_ID as PublicKey.
2. lib/anchor.ts: Load the IDL (copied from /programs/target/idl/agentic_hardware_bridge.json into /dashboard/lib/idl/), construct an anchor Program instance with a read-only provider.
3. hooks/useProgramEvents.ts: Custom hook accepting an event name. Uses program.addEventListener internally. Returns an array of typed events. Handles: reconnect on websocket drop (exponential backoff), dedup by tx signature, cleanup on unmount (program.removeEventListener), max buffer of 100 events.
4. Update PaymentTicker to consume useProgramEvents("ComputePaymentEvent") and ComplianceTable to consume useProgramEvents("ComplianceEvent"). New events animate in from the top.
5. Update TelemetryCard to subscribe to a websocket (NEXT_PUBLIC_BRIDGE_WS_URL) served by the Python bridge process --- this carries off-chain telemetry too high-frequency for Solana.
6. Update TwinViewport to subscribe to ws://localhost:8765 (the PyBullet frame stream) when AUXIN_SOURCE=twin, otherwise render from joint_states pushed via the bridge ws.
7. Connection status indicator in the Header reflects the actual ws state, not a static "Live".
\end{promptbox}

\section{Phase 3: Bridge Service \& Integration}

\subsection*{3.1--3.5: The Bridge Process}
\textbf{Tool:} Cursor IDE \quad \textbf{Model:} Claude Sonnet 4.6

\begin{promptbox}
Context: The bridge is the long-running async Python process that ties everything together: telemetry source -> Gemini oracle -> Solana submission -> dashboard websocket. It must run identically against MockSource, TwinSource, and ROS2Source --- selected only by the AUXIN_SOURCE env var.

Task: Build /sdk/src/auxin_sdk/bridge.py and a /sdk/scripts/run_bridge.py entrypoint.

Requirements:
1. Bridge class. Constructor takes: source (TelemetrySource), oracle (SafetyOracle), program_client (AuxinProgramClient), wallet (HardwareWallet), ws_broadcaster.
2. async def run() main loop: async for frame in source.stream(): await self.process(frame). Cancellable.
3. process(frame) logic:
   a. Push frame to ws_broadcaster (dashboard live telemetry).
   b. If frame has anomaly_flags: hash the frame, await program_client.log_compliance(...), broadcast result. Compliance is unconditional --- never gated on rate limits or budget.
   c. Else: sample a workspace image, await oracle.check(frame, image). If approved: await program_client.stream_payment(small_amount). If denied: hash + log_compliance with severity 1.
4. Submission layer wrapper around program_client: priority fee calculation (use Helius getPriorityFeeEstimate or fall back to a fixed micro-lamport value), blockhash refresh per tx, retry on BlockhashNotFound (max 3), idempotency keys to avoid double-submit on retry.
5. Backpressure: if the submission queue exceeds 50 pending, drop frames whose only flag is informational. NEVER drop compliance events --- they have a separate unbounded queue.
6. WebsocketBroadcaster class: aiohttp ws server on port 8766. Multiple dashboard clients can subscribe. Ping/pong keepalive.
7. /healthz on port 8767 (aiohttp): JSON with source_status, last_successful_tx (signature + slot), last_oracle_latency_ms, queue_depths, uptime.
8. run_bridge.py entrypoint: reads AUXIN_SOURCE, instantiates the right TelemetrySource, wires everything, runs until SIGINT. Structured logs throughout.
9. End-to-end test in tests/test_bridge_e2e.py: spin up Bridge with MockSource against Solana Devnet (skipped in CI without DEVNET_KEYPAIR set), inject an anomaly, assert the matching ComplianceEvent appears on-chain within 5s.
\end{promptbox}

\section{Phase 4: Hardening, Observability, Submission}

\subsection*{4.1--4.3: Observability \& Reproducibility}
\textbf{Tool:} Cursor IDE \quad \textbf{Model:} Claude Sonnet 4.6

\begin{promptbox}
Context: The full Auxin Automata stack works end-to-end on the developer machine. Now we make it reproducible for judges and observable for ourselves.

Task:
1. Prometheus metrics: instrument the bridge with prometheus_client. Counters: auxin_tx_submitted_total{kind, status}, auxin_anomalies_total. Histograms: auxin_oracle_latency_seconds, auxin_solana_submit_latency_seconds. Gauge: auxin_queue_depth{queue}. Expose on :9090.

2. Sentry: add sentry_sdk init in run_bridge.py and Sentry React in /dashboard/app/layout.tsx. SENTRY_DSN from env, both optional.

3. /grafana/auxin_dashboard.json: a checked-in Grafana dashboard JSON with panels for tx rate, oracle latency p50/p95, anomaly count, queue depth.

4. docker-compose.demo.yml at repo root spinning up: bridge (twin mode), dashboard (Vercel build served via `next start`), prometheus, grafana, the twin's ws server. `make demo` runs `docker-compose -f docker-compose.demo.yml up --build` and prints the dashboard URL when ready. Cold start to all-services-healthy in <60s.

5. /scripts/healthcheck.sh that hits /healthz on the bridge and curls a sample tx --- used by docker-compose healthcheck blocks.
\end{promptbox}

\subsection*{4.4--4.5: Documentation}
\textbf{Tool:} Cursor IDE \quad \textbf{Model:} Claude Sonnet 4.6

\begin{promptbox}
Context: Writing the README and architecture docs Colosseum judges will read.

Task: Generate /README.md and per-workspace READMEs.

Root README requirements:
1. One-paragraph elevator pitch (Auxin Automata = autonomous hardware wallets + M2M payments + immutable compliance on Solana).
2. Mermaid architecture diagram showing: Hardware/Twin -> auxin-sdk -> [Gemini Oracle, Anchor Programs on Devnet] -> Dashboard. Arrows labelled with data types.
3. Quickstart: `make demo` -> open localhost:3000. Should work in under 60s on a clean clone with only Docker installed.
4. Deployed addresses table: program ID, sample agent pubkey, sample provider, with Solana Explorer links.
5. Environment variables reference table.
6. Repo layout tree (one-line per workspace).
7. Troubleshooting: top 5 issues and fixes.
8. Team section: Edwin + Tara, links, mention of Aegis at HackEurope and Superteam Ireland outreach.
9. License (Apache 2.0) and contribution note.

Per-workspace READMEs (/sdk, /programs, /edge, /dashboard, /twin): each gets purpose, install, test, run, and how it fits into the bigger picture (with a link back to root).

Tone: technical, confident, no marketing fluff. Judges are engineers and investors --- write for both.
\end{promptbox}

\newpage
\section{Phase 5: Privacy Side Track Integrations}
\textit{These prompts add privacy-preserving M2M payments to the bridge, targeting three Colosseum Frontier side tracks simultaneously. The pattern mirrors the existing AUXIN\_SOURCE agnosticism: a single AUXIN\_PRIVACY env var selects between providers. Build in order: Cloak first (0 competition, earliest deadline May 14), MagicBlock second (deadline May 27), Umbra third (deadline May 26).}

\subsection*{5.0: Privacy Provider Abstraction}
\textbf{Tool:} Claude Code \quad \textbf{Model:} Claude Sonnet 4.6

\begin{promptbox}
Context: We are extending Auxin Automata's bridge to support private M2M payments for three Colosseum Frontier side tracks (Cloak, MagicBlock, Umbra). The current bridge calls program_client.stream_payment() directly as a public SOL transfer. We need to add an abstraction layer so the payment path can be swapped via an env var, exactly like AUXIN_SOURCE controls the telemetry source.

Task: Create the PrivacyProvider abstraction in auxin_sdk/privacy/base.py, update the bridge to use it, and add a DirectProvider that preserves current behaviour.

Requirements:
1. auxin_sdk/privacy/base.py: PrivacyProvider ABC with:
   - async def send_payment(wallet, provider_pubkey, lamports) -> PaymentResult
   - PaymentResult is a Pydantic model: tx_signature (str | None), privacy_provider (str), is_private (bool), confirmation_slot (int | None), metadata (dict)

2. auxin_sdk/privacy/direct.py: DirectProvider(PrivacyProvider) that wraps the existing program_client.stream_payment() call. is_private=False, privacy_provider="direct". This is the current behaviour, preserved exactly.

3. Update bridge.py process() method: instead of calling program_client.stream_payment() directly, call self.privacy_provider.send_payment(). The privacy_provider is injected into the Bridge constructor.

4. Update run_bridge.py entrypoint: read AUXIN_PRIVACY env var (default "direct"). For now, only "direct" is implemented. Add a provider_factory function that maps the string to the concrete class. Raise a clear error for unknown values.

5. Compliance logging is NEVER routed through the privacy provider. Compliance events always go direct to the public chain via program_client.log_compliance(). This is intentional and must be documented in a code comment explaining why.

6. Tests: verify DirectProvider produces identical results to the old direct call. Verify compliance events bypass the privacy provider.

Do NOT implement Cloak, MagicBlock, or Umbra yet. Just the abstraction and the direct passthrough. This is the foundation the next three prompts build on.
\end{promptbox}

\subsection*{5.1: Cloak Integration (Stealth Addresses)}
\textbf{Tool:} Claude Code \quad \textbf{Model:} Claude Opus 4.6 (stealth address crypto needs deeper reasoning)

\begin{promptbox}
Context: Integrating Cloak's stealth address protocol into Auxin Automata for the Cloak side track (deadline May 14, 0 submissions, $5,010 prize pool). Cloak uses stealth meta-addresses: the sender derives a unique one-time address only the recipient can detect and control. A relayer network submits transactions so the stealth address never touches the sender's main wallet.

Cloak docs: https://www.cloakonchain.com/
Cloak is on Solana, uses ZK proofs and Merkle tree commitments for unlinkable withdrawals. Has viewing keys for selective compliance disclosure.

Task: Implement auxin_sdk/privacy/cloak.py as a CloakProvider(PrivacyProvider).

Requirements:
1. Install Cloak's SDK (check their npm/PyPI package --- if TypeScript only, use a subprocess bridge or httpx calls to their API).

2. CloakProvider class:
   - Constructor: takes cloak_api_url (from env CLOAK_API_URL), hardware_wallet.
   - send_payment(): derive a one-time stealth address for the compute provider using Cloak's protocol, send lamports to that stealth address via the relayer, return PaymentResult with is_private=True, privacy_provider="cloak", tx_signature from the relayer response, metadata including the stealth address (for our records, not published).

3. The compute provider must be pre-registered with a stealth meta-address. Add a setup script /scripts/setup_cloak_provider.py that registers the demo provider's stealth meta-address with Cloak. Document this as a one-time setup step.

4. Viewing key integration: store the hardware wallet's Cloak viewing key in the wallet config. This allows authorised auditors to verify payment history without the operator exposing it publicly. Document this in relation to the compliance story: "compliance hashes are public on-chain evidence; payment details are private but auditable via viewing key."

5. Update run_bridge.py: when AUXIN_PRIVACY=cloak, instantiate CloakProvider.

6. Update the dashboard PaymentTicker: when a payment has is_private=True, show a lock icon and "Private Payment" instead of the amount and provider pubkey. Show the tx confirmation but not the stealth address.

7. Error handling: if Cloak's relayer is down or the stealth address derivation fails, fall back to DirectProvider with a warning log. The demo must never stall on a privacy provider failure.

8. Tests: mock Cloak's API, verify stealth address derivation produces unique addresses per payment, verify fallback to direct on failure, verify compliance events are never routed through Cloak.

9. Add /docs/privacy-cloak.md explaining: what Cloak is, why autonomous hardware needs stealth addresses (operational frequency and provider relationships are competitive intelligence), how to enable it, how viewing keys work for auditors, and how this fits the Auxin compliance architecture.
\end{promptbox}

\subsection*{5.2: MagicBlock Integration (Private Ephemeral Rollups)}
\textbf{Tool:} Claude Code \quad \textbf{Model:} Claude Sonnet 4.6

\begin{promptbox}
Context: Integrating MagicBlock's Private Payments API into Auxin Automata for the MagicBlock privacy side track (deadline May 27, 2 submissions, $5,000 prize pool). MagicBlock uses Private Ephemeral Rollups inside TEEs. Users delegate funds, a crank settles them back to Solana with no traceable link between sender and receiver. Compliance is handled at the API layer with wallet screening and AML checks.

MagicBlock Private Payments API docs: check https://docs.magicblock.gg or their GitHub for the REST API spec.

Task: Implement auxin_sdk/privacy/magicblock.py as a MagicBlockProvider(PrivacyProvider).

Requirements:
1. MagicBlockProvider class:
   - Constructor: takes api_url (from env MAGICBLOCK_API_URL), api_key (from env MAGICBLOCK_API_KEY), hardware_wallet.
   - send_payment(): call MagicBlock's REST API to delegate a micro-amount from the hardware wallet into the private ephemeral rollup, specifying the compute provider as the recipient. The rollup's crank settles it. Return PaymentResult with is_private=True, privacy_provider="magicblock", confirmation from the API response.

2. Budget pre-delegation: add a setup function that delegates a compute budget (e.g. 0.1 SOL) into MagicBlock's rollup once, then individual payments draw down from that pool. This avoids a separate delegation per micro-payment. Document the trade-off: pre-delegation is faster but means funds are in the TEE until settled.

3. AML compliance angle: document that MagicBlock runs wallet screening and AML checks on every payment request by default. This is a feature for the submission narrative: "every autonomous M2M payment is AML-screened without any additional infrastructure on the operator's side."

4. Update run_bridge.py: when AUXIN_PRIVACY=magicblock, instantiate MagicBlockProvider.

5. Dashboard: same private payment indicator as Cloak (lock icon, "Private Payment").

6. Fallback to DirectProvider on API failure, with warning log.

7. Tests: mock the MagicBlock REST API, verify delegation and payment flow, verify fallback on failure.

8. Add /docs/privacy-magicblock.md explaining: what MagicBlock's private payments are, why TEE-based settlement suits high-frequency M2M payments (sub-second settlement, no on-chain footprint per individual payment), the AML-by-default angle, and how to enable it.
\end{promptbox}

\subsection*{5.3: Umbra Integration (Shielded Pool)}
\textbf{Tool:} Claude Code \quad \textbf{Model:} Claude Opus 4.6 (ZK/MPC integration needs deeper reasoning)

\begin{promptbox}
Context: Integrating Umbra's privacy protocol into Auxin Automata for the Umbra side track (deadline May 26, 7 submissions, $10,000 prize pool --- largest prize, most competition). Umbra is a privacy layer on Solana using Arcium's MPC network. The TypeScript SDK supports confidential transfers via a unified mixer pool (Merkle trees + ZK proofs), breaking the on-chain link between deposits and withdrawals. Selective disclosure lets users grant viewing access to auditors.

Umbra SDK: TypeScript, available on npm (@umbra/sdk or similar --- check their docs). Supports mainnet, devnet, browser and Node.js.
Umbra docs: https://docs.umbraprivacy.com

Task: Implement auxin_sdk/privacy/umbra.py as an UmbraProvider(PrivacyProvider).

Requirements:
1. Since Umbra's SDK is TypeScript and our bridge is Python, create a lightweight Node.js microservice at /services/umbra-bridge/:
   - Express server with two endpoints: POST /deposit (deposit SOL into Umbra's shielded pool) and POST /withdraw (withdraw to a recipient from the pool).
   - Wraps the Umbra TypeScript SDK calls.
   - Runs as a sidecar process, started by docker-compose alongside the bridge.
   - Communicates with the Python bridge via HTTP on localhost.

2. UmbraProvider class:
   - Constructor: takes umbra_service_url (default http://localhost:3001), hardware_wallet.
   - send_payment(): call the sidecar's /deposit endpoint with the payment amount, then /withdraw to the provider's Umbra address. Return PaymentResult with is_private=True, privacy_provider="umbra", the withdrawal proof hash in metadata.

3. Selective disclosure setup: the hardware wallet's Umbra viewing key is generated during setup and stored alongside the hardware keypair. Add a script /scripts/setup_umbra_viewing_key.py that generates and exports the viewing key. Document: "operators can share this viewing key with auditors or regulators to prove payment history without making it public."

4. Update run_bridge.py: when AUXIN_PRIVACY=umbra, instantiate UmbraProvider. Verify the sidecar is running (health check) before starting the bridge loop.

5. Update docker-compose.demo.yml: add the umbra-bridge sidecar service. Only starts when AUXIN_PRIVACY=umbra.

6. Dashboard: same private payment indicator.

7. Fallback to DirectProvider if the sidecar is unreachable, with warning log.

8. Tests: mock the sidecar HTTP endpoints, verify deposit+withdraw flow, verify fallback on sidecar failure, verify selective disclosure key generation.

9. Add /docs/privacy-umbra.md explaining: what Umbra is, how the shielded pool works conceptually (without leaking their cryptographic internals), why M2M payments benefit from pool-based privacy (individual robot payments are indistinguishable from each other in the pool), the selective disclosure story for auditors, and the sidecar architecture decision.
\end{promptbox}

\subsection*{5.4: Privacy Integration Documentation \& Side Track Submissions}
\textbf{Tool:} Claude Code \quad \textbf{Model:} Claude Sonnet 4.6

\begin{promptbox}
Context: All three privacy providers (Cloak, MagicBlock, Umbra) are implemented. Now we need unified documentation, an updated architecture diagram, and the submission materials for each side track.

Task: Create the unified privacy docs and prepare the side track submission texts.

Requirements:
1. Update the root README.md:
   - Add a "Privacy Providers" section after the existing architecture description.
   - Explain AUXIN_PRIVACY env var: direct (default, public payments), cloak (stealth addresses), magicblock (private ephemeral rollups), umbra (shielded pool).
   - Single table showing: provider name, privacy model, compliance mechanism, setup required.
   - Updated Mermaid architecture diagram adding the privacy provider layer between the bridge and Solana.

2. Create /docs/privacy-overview.md:
   - "Why Autonomous Hardware Needs Payment Privacy": 3 paragraphs covering operational intelligence leakage, competitive risks of public payment patterns, and the compliance paradox (compliance logs must be public, payments should be private).
   - Comparison table of the three providers: Cloak (identity privacy via stealth addresses), MagicBlock (settlement privacy via TEE rollups, AML-by-default), Umbra (transaction privacy via shielded pools, selective disclosure).
   - "Compliance Architecture": explain that compliance events ALWAYS go to the public chain regardless of privacy provider, because immutable public evidence is the entire point of the compliance log. Privacy applies only to the economic layer, never the safety layer.

3. Create three submission text files for the side tracks:
   /docs/submissions/cloak-submission.md
   /docs/submissions/magicblock-submission.md
   /docs/submissions/umbra-submission.md
   Each contains: project name, one-liner tailored to that track's focus, link to repo, link to Devnet program, description of the specific integration (what you built with their SDK, what problem it solves, why M2M hardware payments are a novel use case for their protocol), and a link to the relevant /docs/privacy-*.md file.

4. Tone for submissions: technical, concise, focused on what's novel. Judges for side tracks are the sponsor's team, not Colosseum judges. They want to see their SDK used in a genuinely new way. Lead with "first M2M hardware payment privacy integration on Solana" and back it with the deployed code.
\end{promptbox}

\end{document}
